\documentclass{amsart}
% For dealing with references we use the comment environment
\usepackage{verbatim}
\newenvironment{reference}{\comment}{\endcomment}
%\newenvironment{reference}{}{}
\newenvironment{slogan}{\comment}{\endcomment}
\newenvironment{history}{\comment}{\endcomment}

% For commutative diagrams we use Xy-pic
\usepackage[all]{xy}

% We use 2cell for 2-commutative diagrams.
\xyoption{2cell}
\UseAllTwocells

% We use multicol for the list of chapters between chapters
\usepackage{multicol}

% This is generally recommended for better output
\usepackage{lmodern}
\usepackage[T1]{fontenc}

% For cross-file-references

% Package for hypertext links:
\usepackage{hyperref}

% For any local file, say "hello.tex" you want to link to please

% Theorem environments.
%
\theoremstyle{plain}
\newtheorem{theorem}[subsection]{Theorem}
\newtheorem{proposition}[subsection]{Proposition}
\newtheorem{lemma}[subsection]{Lemma}

\theoremstyle{definition}
\newtheorem{definition}[subsection]{Definition}
\newtheorem{example}[subsection]{Example}
\newtheorem{exercise}[subsection]{Exercise}
\newtheorem{situation}[subsection]{Situation}

\theoremstyle{remark}
\newtheorem{remark}[subsection]{Remark}
\newtheorem{remarks}[subsection]{Remarks}

\numberwithin{equation}{subsection}

% Macros
%
\def\lim{\mathop{\mathrm{lim}}\nolimits}
\def\colim{\mathop{\mathrm{colim}}\nolimits}
\def\Spec{\mathop{\mathrm{Spec}}}
\def\Hom{\mathop{\mathrm{Hom}}\nolimits}
\def\Ext{\mathop{\mathrm{Ext}}\nolimits}
\def\SheafHom{\mathop{\mathcal{H}\!\mathit{om}}\nolimits}
\def\SheafExt{\mathop{\mathcal{E}\!\mathit{xt}}\nolimits}
\def\Sch{\mathit{Sch}}
\def\Mor{\mathop{\mathrm{Mor}}\nolimits}
\def\Ob{\mathop{\mathrm{Ob}}\nolimits}
\def\Sh{\mathop{\mathit{Sh}}\nolimits}
\def\NL{\mathop{N\!L}\nolimits}
\def\CH{\mathop{\mathrm{CH}}\nolimits}
\def\proetale{{pro\text{-}\acute{e}tale}}
\def\etale{{\acute{e}tale}}
\def\QCoh{\mathit{QCoh}}
\def\Ker{\mathop{\mathrm{Ker}}}
\def\Im{\mathop{\mathrm{Im}}}
\def\Coker{\mathop{\mathrm{Coker}}}
\def\Coim{\mathop{\mathrm{Coim}}}

% Boxtimes
%
\DeclareMathSymbol{\boxtimes}{\mathbin}{AMSa}{"02}

%
% Macros for moduli stacks/spaces
%
\def\QCohstack{\mathcal{QC}\!\mathit{oh}}
\def\Cohstack{\mathcal{C}\!\mathit{oh}}
\def\Spacesstack{\mathcal{S}\!\mathit{paces}}
\def\Quotfunctor{\mathrm{Quot}}
\def\Hilbfunctor{\mathrm{Hilb}}
\def\Curvesstack{\mathcal{C}\!\mathit{urves}}
\def\Polarizedstack{\mathcal{P}\!\mathit{olarized}}
\def\Complexesstack{\mathcal{C}\!\mathit{omplexes}}
% \Pic is the operator that assigns to X its picard group, usage \Pic(X)
% \Picardstack_{X/B} denotes the Picard stack of X over B
% \Picardfunctor_{X/B} denotes the Picard functor of X over B
\def\Pic{\mathop{\mathrm{Pic}}\nolimits}
\def\Picardstack{\mathcal{P}\!\mathit{ic}}
\def\Picardfunctor{\mathrm{Pic}}
\def\Deformationcategory{\mathcal{D}\!\mathit{ef}}

\setlength{\emergencystretch}{3em}
\begin{document}
\title[Stacks Project, Tag 01XJ]{363 --- Quasi-coherence and uniform universal vanishing of higher direct images for qcqs morphisms}
\maketitle
\noindent Copyright (C) 2005--2025 Johan de Jong. Stacks Project authors. Source: \href{https://stacks.math.columbia.edu/tag/01XJ}{Tag 01XJ}.

\noindent Editorial difficulty note (not author text). Difficulty H1: The proof must handle nonseparated sources by gluing higher direct images through relative Mayer-Vietoris, not just compute one affine-cover Cech complex. It also constructs a cohomological-degree bound from finite covers and all their intersections that remains valid after every base change..

\noindent Editorial provenance note (not author text). History: excerpt from revision \texttt{a0444\allowbreak 6e57e\allowbreak c1fbc\allowbreak 252a8\allowbreak 71afc\allowbreak ec775\allowbreak 2fb28\allowbreak 07b14}, coherent.tex, lines 741--841. Reference presentation was adapted; original-author context and core support blocks were retained or added. The mathematical wording is unchanged. Licensed under GNU FDL 1.2 or later; no invariant sections or cover texts. See ../licenses/COPYING. Context blocks reproduce author definitions and supporting results; they are not additional counted problems.

\begin{lemma}
\label{lemma-quasi-coherence-higher-direct-images}
Let $f : X \to S$ be a morphism of schemes.
Assume that $f$ is quasi-separated and quasi-compact.
\begin{enumerate}
\item For any quasi-coherent $\mathcal{O}_X$-module $\mathcal{F}$ the
higher direct images $R^pf_*\mathcal{F}$ are quasi-coherent on $S$.
\item If $S$ is quasi-compact, there exists an integer $n = n(X, S, f)$
such that $R^pf_*\mathcal{F} = 0$ for all $p \geq n$ and any
quasi-coherent sheaf $\mathcal{F}$ on $X$.
\item In fact, if $S$ is quasi-compact we can find $n = n(X, S, f)$
such that for every
morphism of schemes $S' \to S$ we have $R^p(f')_*\mathcal{F}' = 0$
for $p \geq n$ and any quasi-coherent sheaf $\mathcal{F}'$
on $X'$. Here $f' : X' = S' \times_S X \to S'$ is the base change of $f$.
\end{enumerate}
\end{lemma}

\begin{proof}
We first prove (1). Note that under the hypotheses of the lemma the sheaf
$R^0f_*\mathcal{F} = f_*\mathcal{F}$ is quasi-coherent by
Schemes, Lemma \href{https://stacks.math.columbia.edu/tag/01LC}{lemma push forward quasi coherent (Tag 01LC)}.
Using
Cohomology, Lemma \href{https://stacks.math.columbia.edu/tag/01E5}{lemma localize higher direct images (Tag 01E5)}
we see that forming higher direct images commutes with restriction
to open subschemes. Since being quasi-coherent is local on $S$ we
reduce to the case discussed in the next paragraph.

\medskip\noindent
Proof of (1) in case $S$ is affine. We will use the induction principle.
Since $f$ quasi-compact and quasi-separated we see that $X$
is quasi-compact and quasi-separated. For $U \subset X$
quasi-compact open and $a = f|_U$ we let $P(U)$ be the property that
$R^pa_*\mathcal{F}$ is quasi-coherent on $S$ for all quasi-coherent modules
$\mathcal{F}$ on $U$ and all $p \geq 0$. Since $P(X)$ is (1), it suffices
the prove conditions (1) and (2) of Lemma \href{https://stacks.math.columbia.edu/tag/08DR}{lemma induction principle (Tag 08DR)} hold.
If $U$ is affine, then $P(U)$ holds because $R^pa_*\mathcal{F} = 0$ for
$p \geq 1$ (by  Lemma \href{https://stacks.math.columbia.edu/tag/01XC}{lemma relative affine vanishing (Tag 01XC)} and
Morphisms, Lemma \href{https://stacks.math.columbia.edu/tag/01SH}{lemma morphism affines affine (Tag 01SH)})
and we've already observed the result holds for $p = 0$ in the first
paragraph. Next, let $U \subset X$ be a quasi-compact open, $V \subset X$
an affine open, and assume $P(U)$, $P(V)$, $P(U \cap V)$ hold.
Let $a = f|_U$, $b = f|_V$, $c = f|_{U \cap V}$, and $g = f|_{U \cup V}$.
Then for any quasi-coherent $\mathcal{O}_{U \cup V}$-module $\mathcal{F}$
we have the relative Mayer-Vietoris sequence
$$
0 \to
g_*\mathcal{F} \to
a_*(\mathcal{F}|_U) \oplus b_*(\mathcal{F}|_V) \to
c_*(\mathcal{F}|_{U \cap V}) \to
R^1g_*\mathcal{F} \to \ldots
$$
see Cohomology, Lemma \href{https://stacks.math.columbia.edu/tag/01EC}{lemma relative mayer vietoris (Tag 01EC)}.
By $P(U)$, $P(V)$, $P(U \cap V)$ we see that
$R^pa_*(\mathcal{F}|_U)$, $R^pb_*(\mathcal{F}|_V)$ and
$R^pc_*(\mathcal{F}|_{U \cap V})$ are all quasi-coherent.
Using the results on quasi-coherent sheaves in
Schemes, Section \href{https://stacks.math.columbia.edu/tag/01LA}{section quasi coherent (Tag 01LA)}
this implies that each of the sheaves
$R^pg_*\mathcal{F}$ is quasi-coherent since it sits in the middle of a short
exact sequence with a cokernel of a map between quasi-coherent sheaves
on the left and a kernel of a map between quasi-coherent sheaves on the right.
Whence $P(U \cup V)$ and the proof of (1) is complete.

\medskip\noindent
Next, we prove (3) and a fortiori (2). Choose a finite affine open
covering $S = \bigcup_{j = 1, \ldots m} S_j$. For each $j$ choose
a finite affine open covering
$f^{-1}(S_j) = \bigcup_{i = 1, \ldots t_j} U_{ji} $.
Let
$$
d_j = \max\nolimits_{I \subset \{1, \ldots, t_j\}}
\left(|I| + t(\bigcap\nolimits_{i \in I} U_{ji})\right)
$$
be the integer found in
Lemma \href{https://stacks.math.columbia.edu/tag/071L}{lemma vanishing nr affines quasi separated (Tag 071L)}.
We claim that $n(X, S, f) = \max d_j$ works.

\medskip\noindent
Namely, let $S' \to S$ be a morphism of schemes and let
$\mathcal{F}'$ be a quasi-coherent sheaf on $X' = S' \times_S X$.
We want to show that $R^pf'_*\mathcal{F}' = 0$ for $p \geq n(X, S, f)$.
Since this question is local on $S'$ we may assume that $S'$ is affine
and maps into $S_j$ for some $j$. Then $X' = S' \times_{S_j} f^{-1}(S_j)$
is covered by the open affines $S' \times_{S_j} U_{ji}$, $i = 1, \ldots t_j$
and the intersections
$$
\bigcap\nolimits_{i \in I} S' \times_{S_j} U_{ji} =
S' \times_{S_j} \bigcap\nolimits_{i \in I} U_{ji}
$$
are covered by the same number of affines as before the base change.
Applying
Lemma \href{https://stacks.math.columbia.edu/tag/071L}{lemma vanishing nr affines quasi separated (Tag 071L)}
we get $H^p(X', \mathcal{F}') = 0$. By the first part of the proof
we already know that each $R^qf'_*\mathcal{F}'$ is quasi-coherent
hence has vanishing higher cohomology groups on our affine scheme $S'$,
thus we see that $H^0(S', R^pf'_*\mathcal{F}') = H^p(X', \mathcal{F}') = 0$
by Cohomology, Lemma \href{https://stacks.math.columbia.edu/tag/01F4}{lemma apply Leray (Tag 01F4)}.
Since $R^pf'_*\mathcal{F}'$ is quasi-coherent
we conclude that $R^pf'_*\mathcal{F}' = 0$.
\end{proof}
\end{document}
